\chapter{camera\_driver\_2}
\label{chap:fpgalix-camera-driver-2}

\section{Introduction}
\label{sec:fpgalix-camera-intro}

This chapter covers how the operating system is instructed to deliver the \term{OV7670} camera's video stream to the applications that need it. The standard solution on \term{Linux} is a kernel driver that uses the \term{V4L2} (\term{Video4Linux2}) framework to expose the camera as a standard device any compatible application can open and read frames from.

|sw/camera_driver_2| is that driver: it captures frames through the \term{Altera} \term{mSGDMA} engine\footnote{See Chapter~\ref{chap:fpgalix-hw-overview}.} into a standard |/dev/videoN| device. Its development was assisted by \term{Claude Code}, an \term{AI} coding assistant, including the kernel-side analysis described in Section~\ref{sec:fpgalix-camera-short-frames}.

\section{V4L2 buffer allocation}
\label{sec:fpgalix-camera-buffer-allocation}

The driver uses \term{videobuf2} (\term{vb2}), the buffer-management layer most \term{V4L2} drivers rely on, which supports several memory models for exchanging frame buffers with \term{userspace}: \term{MMAP} (buffers allocated by the driver, mapped into the application's address space), \term{USERPTR} (buffers allocated by the application, their address passed to the driver instead), and \term{DMABUF} (buffers imported from another device).

Of these, it uses \term{MMAP}, backed by \term{vb2}'s \term{DMA}-contiguous allocator (|vb2_dma_contig_memops|): frame buffers are allocated by the driver itself, each one internally contiguous in physical memory, and \term{userspace} only maps them for reading, through |mmap()|, rather than allocating or owning them. This follows directly from a hardware constraint rather than a design preference: the \term{mSGDMA} engine writes to a single physical start address and length per transfer, so the buffer it writes each individual frame into must itself be physically contiguous, a guarantee arbitrary \term{userspace}-allocated memory cannot provide; the several buffers making up the capture ring need not be contiguous with one another.

\section{Detecting short frames}
\label{sec:fpgalix-camera-short-frames}

The \term{Avalon-ST} pipeline feeding the \term{mSGDMA} marks the end of every frame with an \term{EOP} (\term{End-Of-Packet}) marker. If \term{EOP} arrives before the buffer \term{mSGDMA} was writing to is full, the result is a short frame, and the \term{Altera} \term{mSGDMA} driver already present in the \term{Linux} kernel (|drivers/dma/altera-msgdma.c|), written independently of this project, gave |camera_driver_2| no way to detect this.

This problem was found empirically, through trial and error: streaming from the \term{VHDL} test pattern generators under |hw/my_vhdl/qsys_components/| showed frames coming out wrong. The root cause was then traced by directing \term{Claude Code} to analyse the driver's source directly, over several analysis passes rather than a single one, before the fault was narrowed down to three compounding problems: the \term{DMA} descriptors were not even configured to react to an early \term{EOP} in the first place; once that was corrected, the completion information the hardware did report for each transfer was still being read and then discarded, rather than passed up to |camera_driver_2|; and the first check used to detect this, once it was passed up, proved unreliable on some \term{mSGDMA} hardware revisions and had to be made more robust. All three were fixed.

With this in place, |camera_driver_2| can tell a complete frame from a short one, and, as intended, logs and discards the latter rather than passing it to \term{userspace} as valid.

Two further problems in the same driver were found and fixed: one caused the \term{DMA} engine to stall permanently after the very first short frame, the other caused a second capture session to start from corrupted internal state left behind by the first. Both are documented, with the exact patch, in \cite[``Kernel patch (mandatory)'']{DRV2-README}.

\textbf{Note:} All five changes are already committed to this project's \term{Linux} kernel fork (|linux-socfpga|), which arrives as a submodule. No manual patching is required.

\section{Current limitations}
\label{sec:fpgalix-camera-limitations}

The driver, as it stands, supports a single configuration: \term{Bayer} mosaic (\term{SRGGB8}), 640$\times$480, 30 \term{fps}, matching the \term{OV7670}'s native sensor output and this project's \term{Avalon-ST} pipeline. Other pixel formats and resolutions are left to a future version.

The \term{OV7670} sensor itself is neither initialized nor reset by the driver: this is done separately, over \term{I\textsuperscript{2}C}, before the driver is loaded or a stream started, using the external utilities described in Section~\ref{sec:fpgalix-camera-usage}.

\section{Building}
\label{sec:fpgalix-camera-building}

\begin{shellcode}
cd sw/camera_driver_2
export CROSS_COMPILE=arm-none-linux-gnueabihf-
make
\end{shellcode}

This produces |bin/fpgalix_camera.ko|: the compiled kernel module for this driver. Unlike an ordinary program, a kernel module has no stable, versioned interface to rely on: it links directly into the internal data structures and configuration of one specific kernel build, so it must be compiled against that exact kernel's source tree (Section~\ref{sec:fpgalix-kernel}) to be loadable by it.

\section{Usage}
\label{sec:fpgalix-camera-usage}

Beyond the \term{DE1-SoC} board itself, this requires an \term{OV7670} camera module and the custom camera adapter connecting it to the board's \term{GPIO\_0} header.

\begin{warningbox}
\textbf{Hardware bug:} do not leave the camera in \term{RESET} or \term{POWER-DOWN} for extended periods.\footnotemark
\end{warningbox}
\footnotetext{See Part~\ref{part:appendices}, Section~\ref{sec:camera-adapter-reset-powerdown}.}

Paragraph~\ref{subsec:fpgalix-camera-transfer} shows one way of getting the files onto the board, over |scp|; any other transfer method works just as well. If transferring over \term{SSH}, though, it has to be as the regular user created in Paragraph~\ref{subsec:fpgalix-login-network}: |root| has no working \term{SSH} login. Every step from Paragraph~\ref{subsec:fpgalix-camera-root} onward needs |root|, since loading a kernel module and accessing hardware registers directly both require it; root access can be reached one of two ways:

\begin{itemize}
    \item from that same \term{SSH} session, with |su| (the regular user's own password does not work here: it is the |root| password set in Subparagraph~\ref{subsec:fpgalix-buildroot-base});
    \item from the serial console instead, already logged in as |root|; in this case |cd| into wherever the files were transferred to first, since the |root| login starts elsewhere.
\end{itemize}

\subsection{Deploying the module and utilities}
\label{subsec:fpgalix-camera-transfer}

Copy the compiled module and the debug scripts used in the following steps onto the board together, over |scp|, into a dedicated |camera_driver_2/| folder in the regular user's home directory:

\begin{shellcode}
cd sw/camera_driver_2
ssh <username>@<board-ip> mkdir -p camera_driver_2
scp -p bin/fpgalix_camera.ko utils/debug/resetCtrl.sh utils/debug/ov7670_ctrlif_set.sh
utils/debug/ov7670_dataif_ctrl.sh <username>@<board-ip>:camera_driver_2/
\end{shellcode}

On the board, |cd| into that folder, then run |su| to switch to |root| for the remaining steps, keeping the same working directory:

\begin{shellcode}
cd camera_driver_2
su
\end{shellcode}

\subsection{Releasing the pclk domain reset}
\label{subsec:fpgalix-camera-root}

\begin{shellcode}
./resetCtrl.sh
\end{shellcode}

Select option 2, ``Deassert reset''. This releases |pclk_reset_controller|'s |sw_release| bit, a prerequisite for the camera to respond over \term{I\textsuperscript{2}C} at all.

\subsection{Configuring the OV7670 over I2C}
\label{subsec:fpgalix-camera-i2c-config}

\begin{shellcode}
./ov7670_ctrlif_set.sh
\end{shellcode}

Select the \term{Bayer} mode (1 for \term{Raw}, 2 for \term{Processed}). Registers are written and then read back automatically, confirming the camera answered correctly.

\subsection{Enabling frame acquisition}
\label{subsec:fpgalix-camera-frame-acq}

\begin{shellcode}
./ov7670_dataif_ctrl.sh
\end{shellcode}

Select option 1, ``Enable frame acquisition''.

\subsection{Loading the driver}
\label{subsec:fpgalix-camera-load}

\begin{shellcode}
insmod fpgalix_camera.ko
\end{shellcode}

\subsection{Viewing the stream}
\label{subsec:fpgalix-camera-viewing}

|utils/debug/view.py| is a small \term{OpenCV} viewer, run on a \term{PC}, that displays the raw frames read from the board's |/dev/video0| over \term{SSH}:

\begin{shellcode}
ssh <username>@<board-ip> "v4l2-ctl -d /dev/video0 --stream-mmap --stream-to=-" | python3 view.py -dwg
\end{shellcode}

|v4l2-ctl|, run on the board, streams raw frames to its own standard output instead of a file; piped over \term{SSH}, |view.py| reads them from its standard input and displays them. Its processing modes (|-r/-d/-dw/-dwg/-F|) are documented in \cite[``view.py --- Live preview'']{DRV2-README}. Its dependencies, |python3|, |python3-numpy| and |python3-opencv|, are already covered in Paragraph~\ref{subsec:fpgalix-apt-packages}.

|v4l2-ctl| itself reads directly from |/dev/video0| and needs no elevated privileges beyond read access to that device node.
